Ce projet de fin d'études avait pour ambition de répondre à une problématique concrète~:
\textit{comment intégrer la sécurité à chaque étape d'une chaîne de livraison logicielle
cloud-native, sans sacrifier la vélocité~?} La conclusion générale revient sur l'ensemble
de la démarche, évalue l'atteinte des objectifs, analyse les difficultés rencontrées,
mesure les apports de part et d'autre, et ouvre sur des perspectives d'approfondissement.

\section*{Récapitulation de la Démarche}

La démarche suivie tout au long de ce projet s'est organisée en quatre étapes progressives
et complémentaires. Un premier travail d'analyse a permis de cartographier précisément
les lacunes du pipeline Jenkins existant~: absence totale de scans de sécurité, clés
d'authentification statiques, infrastructure non reproductible et absence de monitoring.
Sur la base de cette analyse, les besoins fonctionnels et non fonctionnels ont été formalisés,
puis une vision cible DevSecOps a été conçue autour de huit principes directeurs et d'une
architecture en six couches.

L'implémentation, organisée en quatre sprints Agile, a couvert~: le provisionnement de
l'infrastructure GCP via Terraform (Sprint~1), la containerisation et le pipeline CI/CD
avec authentification keyless (Sprint~2), l'intégration des couches de sécurité SAST/SCA/DAST
et le déploiement GitOps avec réseau Zero Trust (Sprint~3), et enfin les tests E2E,
l'observabilité et la surveillance runtime (Sprint~4).

\section*{Présentation des Résultats}

Les résultats obtenus répondent directement à chaque lacune identifiée dans la problématique
initiale. Le tableau~\ref{tab:resultats_finaux} synthétise les indicateurs clés.

\begin{table}[h]
\centering
\caption{Bilan quantitatif de la transformation DevSecOps}
\label{tab:resultats_finaux}
\begin{tabularx}{\textwidth}{|l|X|X|}
\hline
\rowcolor{gray!20}\textbf{Indicateur} & \textbf{Avant} & \textbf{Après} \\
\hline
CVEs détectées (npm) & 0 (non scanné) & 9 CVEs dont 4 HIGH \\
\hline
CVEs détectées (image) & 0 (non scanné) & 47 CVEs dont 3 CRITICAL \\
\hline
Findings OWASP ZAP & 0 (non testé) & 10 findings dont 1 HIGH (XSS) \\
\hline
Secrets détectés & 0 (non scanné) & 3 secrets exposés identifiés \\
\hline
Auth CI/CD & Clé JSON statique & OIDC éphémère (WIF) \\
\hline
Infrastructure & ClickOps manuel & Terraform IaC 100\% \\
\hline
Tests E2E & Absents & 6 scénarios Playwright in-cluster \\
\hline
Coverage tests unitaires & Non mesurée & 70\,\%+ (Vitest) \\
\hline
NetworkPolicies & Aucune & 12 règles default-deny \\
\hline
SIEM actif & Absent & Wazuh 4 agents opérationnels \\
\hline
ZAP FAIL count & Non applicable & 0 (après durcissement nginx) \\
\hline
\end{tabularx}
\end{table}

\section*{Difficultés Rencontrées}

\textbf{Contraintes GKE Autopilot~:} Le mode Autopilot bloque les capacités
privilégiées des pods (\texttt{hostPID}, \texttt{hostNetwork}). Ces restrictions ont
limité la surveillance kernel de Wazuh. La solution retenue est de déployer le manager
Wazuh sur une VM Compute Engine dédiée et de connecter les agents via TCP~1514.

\textbf{Orchestration des tests E2E en GitOps~:} L'intégration de Playwright comme
PostSync hooks Argo CD a nécessité une conception minutieuse des RBAC et des
NetworkPolicies pour permettre aux pods de tests d'accéder à l'application dans le
respect du principe Zero Trust.

\textbf{Gestion des faux positifs~:} Les scanners génèrent des faux positifs qui
nécessitent un tri manuel. La mise en place de fichiers d'exception (\texttt{.trivyignore},
règles SonarCloud) a permis de filtrer le bruit sans masquer les vulnérabilités réelles.

\textbf{Workload Identity Federation~:} La configuration initiale du mécanisme OIDC
entre GitHub Actions et GCP a exigé une maîtrise des concepts d'identité fédérée,
de pools de fournisseurs d'identité et de conditions d'attribut~---~concepts peu documentés
pour ce cas d'usage précis.

\section*{Apports du Stage}

\textbf{Apports à l'entreprise~:} Le projet a livré à Sopra Steria une plateforme
DevSecOps opérationnelle et documentée, un référentiel de bonnes pratiques, une
démonstration concrète de l'intégration SAST/SCA/DAST dans GitHub Actions, et des
templates Terraform/Kubernetes réutilisables pour d'autres projets.

\textbf{Apports pour l'élève-ingénieur~:} Ce stage a permis d'acquérir des compétences
transversales rares combinant le cloud engineering, la sécurité applicative, l'orchestration
Kubernetes et l'ingénierie des systèmes distribués.

\begin{table}[h]
\centering
\caption{Compétences acquises durant le stage}
\label{tab:competences}
\begin{tabularx}{\textwidth}{|l|X|}
\hline
\rowcolor{gray!20}\textbf{Domaine} & \textbf{Compétences Acquises} \\
\hline
Cloud Engineering & GCP, GKE Autopilot, Terraform, Artifact Registry, IAM \\
\hline
DevSecOps & GitHub Actions, Argo CD, GitOps, Shift-Left Security \\
\hline
Sécurité applicative & Trivy, OWASP ZAP, SonarCloud, gestion des CVEs \\
\hline
Kubernetes & Deployments, NetworkPolicies, RBAC, HPA, Jobs, Ingress \\
\hline
Observabilité & Prometheus, Grafana, ServiceMonitor, métriques nginx \\
\hline
Développement & React 18, TypeScript, Vite, Vitest, Playwright E2E \\
\hline
Méthode & Analyse de risque, architecture en couches, documentation \\
\hline
\end{tabularx}
\end{table}

\section*{Perspectives d'Approfondissement}

\textbf{Court terme~:}
\begin{itemize}
  \item Intégration de \textbf{Gitleaks} en hook pre-commit pour détecter les secrets
        avant qu'ils n'atteignent le dépôt Git~;
  \item Signature des images Docker avec \textbf{cosign} et vérification à l'admission
        via \textbf{Kyverno}~;
  \item Ajout de règles Kyverno pour enforcer les contraintes de sécurité pod-level.
\end{itemize}

\textbf{Moyen terme~:}
\begin{itemize}
  \item Migration vers \textbf{Istio Service Mesh} pour une mTLS automatique entre
        services et une observabilité réseau L7~;
  \item Génération et publication d'un \textbf{SBOM} (\textit{Software Bill of Materials})
        via Syft à chaque release~;
  \item Intégration de \textbf{Falco} pour la détection d'anomalies comportementales
        syscall dans les pods.
\end{itemize}

\textbf{Long terme~:}
\begin{itemize}
  \item Atteindre le niveau \textbf{SLSA~3} avec des preuves de provenance vérifiables~;
  \item Déploiement multi-région avec stratégie de \textit{disaster recovery} automatisée~;
  \item Mise en conformité formelle \textbf{ISO~27001} en s'appuyant sur les artefacts
        d'audit générés automatiquement par la plateforme.
\end{itemize}

\section*{Message Final}

Ce projet illustre une vérité fondamentale du développement logiciel moderne~:

\begin{center}
\large\textit{«~La sécurité n'est pas un obstacle à la vélocité~--- c'est un enabler de confiance.~»}
\end{center}

Une équipe qui automatise ses contrôles de sécurité livre plus vite, avec plus de
confiance, et consacre moins de temps à gérer des incidents. Le DevSecOps n'est pas
une contrainte imposée par les équipes sécurité~--- c'est une pratique d'ingénierie
mature qui bénéficie à l'ensemble de l'organisation.
